\chapter{Autopsy of a Prototype}\label{ch:autopsy}

\begin{quote}\small\itshape
The project did not begin with a blank page. It began with a working ---
running, loss-decreasing, plausible-looking --- prototype that was wrong in
eight distinct ways, none of which produced an error message. This chapter
dissects it. The point is not self-flagellation: the prototype's
\emph{architectural idea} was sound, and the corrected system keeps it. The
point is that each defect is a specimen of an entire \emph{class} of failure,
and the corrected system is best understood as eight class-level defenses.
Every later chapter of this book repairs something that dies on this table.
\end{quote}

\section{The specimen}\label{sec:specimen}

The prototype (preserved untouched in the repository at
\code{OP model/nec/nec\_hybrid\_explained.py}) implements a
\emph{normal-versus-extreme conditional} model --- NEC --- in three pieces:

\begin{itemize}[leftmargin=2em]
\item \code{C\_L}: a stacked LSTM sequence encoder (defaults: hidden size
  $256$, \emph{eight} layers) reading the feature window;
\item \code{C\_F}: a classifier head on one selected time step of the LSTM
  output --- \code{Linear} $\to$ \code{BatchNorm1d} $\to$ activation $\to$
  \code{Linear} $\to$ \code{nn.Softmax} --- emitting probabilities over
  \emph{three} market classes;
\item two \code{NEC\_MLP} regressors, ``\code{extreme}'' and
  ``\code{normal},'' each reading a raw slice of the input window.
\end{itemize}

The forward pass routes hard:

\begin{lstlisting}
prob  = self.classifier(h)               # (B, 3) softmax output
e_out = self.extreme(x[:, -1, :-1])      # slice: last step, drop last col
n_out = self.normal(x[:, -1, :-1])
pred  = torch.where(torch.argmax(prob, dim=1) == 1, e_out, n_out)
\end{lstlisting}

Figure~\ref{fig:prototype} draws the whole machine with the defects marked
where they live.

\begin{figure}[t]
\centering
\begin{tikzpicture}[
  font=\small,
  box/.style={draw, rounded corners=1.5pt, minimum height=8mm,
              minimum width=26mm, align=center, fill=white},
  bad/.style={draw=breakframe, circle, fill=breakbg, inner sep=1pt,
              font=\bfseries\footnotesize, text=breakframe},
  arr/.style={-{Latex[length=2mm]}, thick},
  ]
  \node[box] (xseq) {window $x$\\ $(B, T, d)$};
  \node[box, right=12mm of xseq] (lstm) {\code{C\_L}: LSTM\\ $n{=}256$, $L{=}8$};
  \node[box, right=12mm of lstm] (cf) {\code{C\_F}: head\\ BatchNorm $\to$ Softmax};
  \node[box, right=12mm of cf] (route) {\code{argmax} $=1$?};
  \node[box, below=13mm of xseq] (slice) {\code{x[:, -1, :-1]}\\ raw last-step slice};
  \node[box, below=13mm of cf] (ext) {\code{extreme} MLP};
  \node[box, below=8mm of ext] (norm) {\code{normal} MLP};
  \node[box, below=13mm of route] (pred) {$\hat y$ (one expert,\\ winner takes all)};

  \draw[arr] (xseq) -- (lstm);
  \draw[arr] (lstm) -- (cf);
  \draw[arr] (cf) -- (route);
  \draw[arr] (xseq) -- (slice);
  \draw[arr] (slice.east) -- (ext.west);
  \draw[arr] (slice.east) -- (norm.west);
  \draw[arr] (ext.east) -- (pred.west);
  \draw[arr] (norm.east) -| (pred.south);
  \draw[arr] (route) -- (pred);

  \node[bad, above right=-1.5mm and -3mm of lstm] {5};
  \node[bad, above right=-1.5mm and -3mm of cf] {3,4};
  \node[bad, above right=-1.5mm and -3mm of route] {1,2};
  \node[bad, below right=-1.5mm and -3mm of slice] {6,7};
  \node[bad, above left=-1.5mm and -2mm of lstm] {8};
\end{tikzpicture}
\caption{\textbf{The prototype, with defects in place.} Numbers key to the
ledger of Section~\ref{sec:ledger}. Note the two load-bearing absurdities
visible even at block-diagram level: the gradient cannot cross the
\code{argmax} (1), so the entire left half of the diagram --- the
four-million-parameter encoder and its head --- receives no training signal
from the prediction task; and the experts read a raw, undocumented slice of
the input (6, 7), so the encoder's summary influences \emph{nothing} except
the untrainable routing decision.}
\label{fig:prototype}
\end{figure}

Before the knife comes out, record what is \emph{right}. The prototype's
thesis --- encode the recent sequence into a state, classify the market
regime from that state, and route to a regime-specialized regressor --- is
exactly Definition~\ref{def:nec}. The corrected system keeps this skeleton
wholesale: encoder $\to$ gate $\to$ experts. Diagnosing which part of a
broken system is sound is as much a skill as finding the breaks; the ledger's
verdict on the architecture was \emph{keep}, and the whole project is
downstream of that judgment call being correct.

\section{The ledger}\label{sec:ledger}

Eight defects, in logical rather than severity order. For each: what the
code does, why it is wrong (with the mathematics where mathematics is the
right scalpel), and the repair --- with a pointer to the chapter that builds
it properly.

\subsection*{Defect 1 --- \code{argmax} routing: the gate cannot learn}

The routing decision \code{torch.where(argmax(prob) == 1, e\_out, n\_out)}
selects one expert per sample. As a function of the classifier's parameters
$\theta$, the selection $\ind{\argmax_k \mathrm{prob}_k(\theta) = 1}$ is
piecewise constant: small perturbations of $\theta$ leave the argmax
unchanged almost everywhere, so
\[
  \frac{\partial\, \hat y}{\partial \theta} = 0
  \quad \text{almost everywhere.}
\]
Autograd dutifully computes this zero. Whatever loss is attached to
$\hat y$, \emph{no gradient ever reaches the classifier} through the routing
path --- the gate's weights receive exactly zero training signal from the
prediction task. The four-million-parameter encoder feeding it is trained,
if at all, only by whatever auxiliary classification loss the gate carries,
and the routing quality can never improve from prediction feedback. Nothing
crashes; the loss on the experts still decreases; the model ``trains.''

The repair (Chapter~5) is to stop making a decision and start mixing:
weight each expert's \emph{density} by a probability and maximize the
mixture likelihood. The gradient of that objective with respect to the gate
logits will turn out to be the beautiful $\pi - r$ --- prior minus posterior
--- restoring a rich training signal to precisely the module that had none.
The lesson generalizes far beyond this file:
\begin{center}
\emph{hard decisions inside a differentiable model are gradient graves;}\\
\emph{make the decision soft, or move it outside the training path.}
\end{center}
And the story has a coda that will recur in Chapter~6: the ``obvious'' soft
fix of renormalizing the top-1 weight is \emph{equally} gradient-dead --- a
constant $1$ is as flat as an indicator --- a fact the corrected package
enforces with a validation error and proves with a test.

\subsection*{Defect 2 --- three classes collapsed to one bit}

The classifier emits three market classes; the routing consumes only
$\ind{\text{class} = 1}$. Classes $0$ and $2$ --- presumably ``calm'' and
some third state --- are silently merged, the distinction the classifier
spends capacity learning is discarded, and the number of regimes is
hard-wired into an equality test buried in the forward pass. $K$ is not a
parameter; it is folklore.

The repair is structural, not cosmetic: \emph{all} $K$ components enter the
likelihood with their own weights and densities (Chapter~5), $K$ lives in
one configuration field, and no class index has hardcoded meaning anywhere
--- which later forces an honest confrontation with the fact that mixture
components are exchangeable (the label-permutation problem, Chapters~7--8),
a problem the prototype never met only because it had welded the labels to
the furniture.

\subsection*{Defect 3 --- \code{nn.Softmax} inside the model, feeding a
second loss}

The head ends in \code{nn.Softmax}, so downstream code sees probabilities,
and any classification loss must be computed as
$\log(\mathrm{softmax}(z))$ in two separate steps. Two distinct failures
hide here. \emph{Numerical:} $\log(\mathrm{softmax}(z))$ computed
na\"ively overflows for logits as small as $z \approx 90$ (since
$e^{90} \approx 10^{39}$ exceeds float32's range) and catastrophically
cancels for confident predictions, where the winning probability rounds to
$1.0$ and its log to exactly $0$; the fused \code{log\_softmax} identity
$\log \pi_k = z_k - \LSE(z)$, with the max-subtracted log-sum-exp, is exact
and stable at any scale --- Chapter~11 makes log-domain computation a
system-wide law. \emph{Statistical:} the gate was to be trained by a
\emph{separate} classification objective against externally supplied
regime labels, while the experts trained on regression --- two losses, two
sources of truth, and a gate whose notion of ``regime'' is whatever the
labeler said rather than whatever best explains the returns. The corrected
model has \emph{one} scalar objective, the mixture NLL, in which routing and
prediction are the same optimization (Chapter~5) and the gate is
unsupervised by design (Decision~B, Chapters~2 and~9).

\subsection*{Defect 4 --- \code{BatchNorm1d} in the head: samples that
whisper to each other}

BatchNorm normalizes each activation by statistics of the \emph{current
batch}. Consequences, in ascending severity for this domain: (i)
train/eval asymmetry --- at eval time the layer switches to running
statistics, so the very function being computed differs between training
and deployment; (ii) batch-composition sensitivity --- a sample's output
depends on which other samples share its batch; (iii) \emph{leakage} ---
in a panel, a batch can mix dates, and normalizing by batch statistics lets
information flow between dates within the batch: sample $A$ at date $t$ is
transformed using moments computed partly from samples at $t' > t$. That is
look-ahead, laundered through a normalization layer, invisible in any code
review that does not know to look for it. And by
Remark~\ref{rem:ratchet}, if it helps the backtest, nothing downstream
will ever flag it.

The repair is the drop-in \code{LayerNorm} --- identical normalization
\emph{within} each sample across features, no cross-sample statistics, no
train/eval asymmetry, no leakage channel (Chapter~11). The general lesson:
in time-ordered data, \emph{any} operation that couples samples --- batch
normalization, batch-level feature scaling, cross-sectional ranks computed
over a batch rather than a date --- is a leakage suspect until proven
innocent, and the proof must be structural (the operation cannot see the
future) rather than empirical (``the backtest looks fine'').

\subsection*{Defect 5 --- four million parameters against a whisper of
signal}

The encoder defaults: LSTM, hidden size $n = 256$, $L = 8$ stacked layers.
Parameter count for an LSTM layer is $4\bigl(n(n + d) + n\bigr)$ --- four
gates, each with input and recurrent weights and a bias --- so with $d$
small, layer one costs $4(256 \cdot 256 + 256 \cdot d + 256) \approx 270$K
and layers two through eight, whose input dimension is $n$, cost
$4(256^2 + 256^2 + 256) \approx 525$K each: roughly \emph{four million
trainable parameters} before the head. Chapter~\ref{ch:cross} priced the
other side of the scale: a generous panel offers order $10^6$ rows whose
effective information content is orders of magnitude lower
(\S\ref{sec:dependence}), pursuing a signal of $R^2 \sim 10^{-3}$. This is
not a subtle bias--variance trade-off to be tuned; it is a category error,
resolvable by arithmetic before any experiment is run.

The repair (Chapters~9 and~11): a one-layer GRU with $n = 32$
--- about $3\bigl(n(n+d)+n\bigr) \approx 3{,}500$ parameters --- with the
GRU-over-LSTM choice itself an application of the same arithmetic (three
gates versus four: a $25\%$ parameter cut at equal hidden size is pure
variance reduction at this data scale). The deeper practice being taught:
\emph{effective-sample arithmetic is a design-time activity}. Parameter
budgets in this field are set by the noise floor, not by what fits in GPU
memory, and ``bigger model'' claims must clear a capacity-matched baseline
(Chapter~16) precisely because capacity itself is a confounder.

\subsection*{Defect 6 --- the encoder's summary influences nothing
trainable}

Follow the information flow in Figure~\ref{fig:prototype}. The experts read
\code{x[:, -1, :-1]} --- a raw slice of the input window --- not the
encoder's state. The encoder state reaches only the classifier, whose
routing decision is gradient-dead by defect~1. Net effect: the LSTM's
summary of history influences no trainable pathway. Four million parameters
compute a representation whose only consumer cannot learn. The model is,
functionally, two MLPs on a raw feature slice plus an elaborate,
untrained switch.

The repair is a genuine \emph{decision}, not a reflex --- because once
defect~1 is fixed and the gate trains, ``what should the experts see?''
has two defensible answers. Feed experts the snapshot features only, and
the encoder specializes in \emph{state} while experts specialize in
\emph{cross-section} --- clean division of labor, fewer parameters, and the
regime interpretation of Chapter~2 stays legible. Feed them snapshot
$\oplus$ encoder state, and experts can condition on history --- more
capacity, blurrier semantics, and the gate/expert division of labor can
collapse (if the experts see the state, they can privately un-mix the
regimes and the gate's job evaporates). The corrected design makes this
Decision~A: snapshot-only by \emph{default}, switchable by configuration,
with the trade-off documented and the switch exposed as an ablation
(Chapter~9). Defects that are really decisions deserve to be decided in
daylight.

\subsection*{Defect 7 --- \code{x[:, -1, :-1]}: the undocumented reserved
column}

The slice quietly asserts a convention: \emph{the last feature column is
not a feature}. Somewhere upstream, the data pipeline packed something ---
by all evidence the \emph{target} --- into the final column of the feature
tensor, and every consumer must remember to slice it off. The contract
lives in a \code{-1} in the middle of a forward pass: undocumented,
untyped, unchecked. Now enumerate the failure modes. A new feature appended
after the reserved column shifts the target \emph{into} the model's view:
instant, silent, catastrophic look-ahead (the model ``predicts'' returns it
was shown). A consumer that forgets to slice trains on the target directly.
A pipeline change that drops the reserved column makes every downstream
slice off by one --- silently, since shapes still match. By
Remark~\ref{rem:ratchet}, the first failure mode is the one you will
never find, because it makes results \emph{better}.

The repair (Chapter~12): a named feature schema --- every column of both
tensors has a name checked at panel construction; the target is a separate
tensor of the \code{Panel}, so there \emph{is no} reserved column to slice
off; and the model never slices feature tensors, anywhere. The class-level
lesson: \emph{data-layout conventions are contracts, and contracts that
live in comments (or worse, in a \code{-1}) eventually get broken by
someone who never saw them.} Encode contracts in types, schemas, and
constructors that refuse to build the invalid object.

\subsection*{Defect 8 --- default dimensions that disagree}

\code{C\_L} defaults its hidden size to $256$; \code{C\_F} defaults its
expected input dimension to $64$. Instantiated with defaults, the pieces do
not fit --- and whether that produces a loud shape error or a silent
mis-projection depends on incidental details of which linear layer touches
the tensor first. The defect is not the mismatch itself but the
\emph{design that made it possible}: each module carries its own copy of a
shared quantity, and nothing owns the truth. Configuration drift is the
generic name --- two defaults, one meaning, no reconciliation.

The repair (Chapters~9 and~11): a single configuration tree in which every
dimension is written exactly once, cross-validated at construction time
(\code{NECConfig.validate()} checks every derived dimension and raises with
names and numbers), with derived quantities \emph{computed}, never
restated --- the gate's input dimension, for instance, is a property
derived from the encoder's hidden size and is deliberately not a settable
field, making the prototype's mismatch \emph{unrepresentable}. Shape
assertions at module boundaries (\code{assert\_shape}, with tensor names in
the error) catch what validation cannot. The slogan for the whole family:
\emph{make illegal states unrepresentable, and make the remaining errors
loud and early.}

\section{A taxonomy, and the shape of the cure}\label{sec:taxonomy}

Eight defects, four families --- and the families, not the instances, are
what the rest of the book is organized around:

\begin{center}
\small
\begin{tabular}{@{}p{0.26\textwidth}cp{0.42\textwidth}c@{}}
\toprule
family & defects & defense layer & built in \\
\midrule
\textbf{Objective pathologies} --- the loss cannot train what it must
  & 1, 2, 3 & derive the objective and its gradients \emph{before} coding;
  fused log-domain losses; gradient-flow regression tests & Part II \\
\addlinespace
\textbf{Leakage channels} --- information crosses the time boundary
  & 4, 7 & data contracts (schema, separate target tensor); leakage stated
  as executable invariants; no cross-sample ops in models & Parts III, IV \\
\addlinespace
\textbf{Capacity delusions} --- parameters vastly outnumber effective
  samples
  & 5, 6 & effective-sample arithmetic at design time; small defaults;
  capacity-matched baselines; decisions-as-ablations & Parts II, III, V \\
\addlinespace
\textbf{Contract failures} --- modules disagree about shared truths
  & 8 & single config tree, validated at construction; derived (not
  settable) dimensions; shape assertions at boundaries & Part III \\
\bottomrule
\end{tabular}
\end{center}

Notice what is \emph{absent} from the defense column: ``be more careful.''
None of the eight defects would have been prevented by vigilance, because
each was invisible at the moment of writing --- the code ran, the loss fell,
and (Remark~\ref{rem:ratchet}) the failures that helped the numbers had
negative incentive to be found. Every defense above is \emph{structural}:
a derivation that must exist before the loss is coded, a constructor that
refuses the invalid object, a test that fails on the defect, an API that
does not offer the dangerous operation. The prototype's deepest lesson is
that in this field, correctness is a property of \emph{systems}, not of
authors.

\begin{engineering}[every defect becomes a test]
The corrected package operationalizes the ledger as a regression suite:
for each defect, a test that \emph{would have failed on the prototype}.
Defect~1: after one backward pass of the mixture NLL, the gate's weight
gradient is asserted nonzero --- and a companion test builds an
argmax-routing double and asserts its gate gradient is \emph{exactly} zero,
pinning the defect itself, not just the fix. Defect~4: LayerNorm's output
for a sample is asserted invariant to its batch companions. Defect~7: a
target column smuggled into a feature tensor changes its width and fails
the schema check. Defect~8: constructing the config with mismatched
dimensions raises at build time. Tests-as-theorems is the working style of
the whole project --- Chapter~19 is devoted to it --- and it began here, as
the only honest response to a prototype that had passed every informal
inspection.
\end{engineering}

\section{The road ahead}\label{sec:roadmap}

The book now proceeds in the only defensible order: theory first, because
five of the eight defects (1, 2, 3, 5, 6) are visible \emph{only} to someone
who has derived the objective and priced the statistics --- no amount of
code review finds a gradient grave if the reviewer does not know what the
gradient must be.

\begin{center}
\small
\begin{tabular}{@{}clp{0.52\textwidth}@{}}
\toprule
defect & repaired in & by \\
\midrule
1 & Ch.~5, 6 & mixture NLL; the $\pi{-}r$ gate gradient; trainable hard
              routing done right (hard-EM), top-1 renormalization proven
              dead \\
2 & Ch.~5, 7, 8 & all-$K$ likelihood; $K$ as configuration; label
              permutation confronted honestly \\
3 & Ch.~5, 11 & one fused objective; log-domain law; logits-only modules \\
4 & Ch.~11, 12 & LayerNorm; no cross-sample operations; leakage invariants \\
5 & Ch.~9, 11 & effective-sample arithmetic; GRU $n{=}32$; capacity-matched
              baselines (Ch.~16) \\
6 & Ch.~9 & Decision~A --- expert inputs decided, documented, switchable \\
7 & Ch.~12 & feature schema; target outside the feature tensors \\
8 & Ch.~9, 11 & single validated config tree; derived dimensions; shape
              assertions \\
\bottomrule
\end{tabular}
\end{center}

One closing observation. The prototype was not written by a fool; it was
written by the same intelligence that later diagnosed all eight defects.
What changed in between was not talent but \emph{method}: derivations
before code, arithmetic before architecture, contracts before conventions,
and tests that encode the failures already survived. That method is the
actual subject of this book. The mixture-of-experts is merely the occasion.

\exercisehead

\begin{exercise}\label{ex:3-argmax}
(\emph{The gradient grave, exhibited.}) In five lines of PyTorch, build
$z \in \R^3$ with \code{requires\_grad=True}, compute
$\hat y = \code{torch.where}(\argmax_k \mathrm{softmax}(z)_k = 1,\, a,\, b)$
for constants $a \ne b$, backpropagate a loss through $\hat y$, and print
\code{z.grad}. Explain the result via the piecewise-constance argument of
defect~1. Then replace the hard selection with the probability-weighted
combination $\hat y = \sum_k \pi_k \mu_k$ and show the gradient comes
alive. (This is, in miniature, the package's
\code{test\_gate\_receives\_gradient} and its argmax companion.)
\end{exercise}

\begin{exercise}\label{ex:3-params}
(\emph{Capacity arithmetic.}) (a) Derive the LSTM per-layer parameter count
$4\bigl(n(n + d_{\text{in}}) + n\bigr)$ from the four-gate equations, and
the GRU's $3\bigl(n(n + d_{\text{in}}) + n\bigr)$. (b) Verify the
$\approx 4$M figure for the prototype ($n = 256$, $L = 8$, small $d$) and
compute the corrected encoder's count ($\text{GRU}, n = 32, L = 1, d = 2$).
(c) Using Chapter~\ref{ch:cross}'s panel ($N = 500$ names, $T = 2{,}500$
dates) and the crude rule ``at most one parameter per ten effective
samples,'' bound the affordable parameter count if effective samples are
(i) all rows, (ii) rows discounted by cross-sectional effective size
$N_{\text{eff}} = 20$. Where does each encoder sit?
\end{exercise}

\begin{exercise}\label{ex:3-bn}
(\emph{Hands-on: BatchNorm leaks, LayerNorm does not.}) Construct a toy
two-date batch where date $2$'s samples have a large common shift
(simulating a future volatility spike). Pass it through
\code{nn.BatchNorm1d} in training mode and show date $1$'s
\emph{normalized} outputs change when date $2$'s samples are included
versus excluded --- information flowed backward in time. Repeat with
\code{nn.LayerNorm} and show date $1$'s outputs are bit-identical in both
cases. State, in one sentence each, which of the three BatchNorm
consequences of defect~4 your experiment exhibited.
\end{exercise}

\begin{exercise}\label{ex:3-reserved}
(\emph{The reserved column, weaponized.}) Simulate defect~7: build a
feature matrix whose final column \emph{is} the target, train (i) a model
that correctly slices it off and (ii) a model that ``forgets,'' and compare
their in-sample and out-of-sample $R^2$ on a task where the true signal is
$R^2 = 0.01$. Then append one extra innocent feature \emph{after} the
reserved column (the shift failure) and measure again. Write the two-line
schema check that would have made all three configurations fail loudly at
construction.
\end{exercise}

\begin{exercise}\label{ex:3-taxonomy}
(\emph{Classify new specimens.}) Assign each of the following real-world
bugs to a family of Section~\ref{sec:taxonomy} and name the structural
defense that catches it: (a) a validation split created with
\code{sklearn.train\_test\_split(shuffle=True)} on a daily panel; (b) a
z-score normalization whose mean and standard deviation were fitted on the
full history before the train/test split; (c) an attention model whose
mask allows position $t$ to attend to $t{+}1$ during training only; (d) a
config file where \code{seq\_len} appears in three places and a refactor
changes two of them; (e) a reinforcement-learning reward computed from
next-day prices already visible in the state. (One sentence each.)
\end{exercise}

\begin{exercise}\label{ex:3-keep}
(\emph{The judgment call.}) The ledger's verdict on the prototype's
architecture was \emph{keep}. Write one paragraph for the opposite brief:
the strongest honest argument that the encoder--gate--experts skeleton
should have been abandoned (consider: Chapter~\ref{ch:cross}'s variance
budget, the optimization pathologies previewed in
\S\ref{sec:whynotfeatures}, and what a plain pooled ridge on snapshot
features costs to build and maintain by comparison). Then write the
rebuttal you find most convincing. Keep both; Part~V will referee.
\end{exercise}
